\documentclass[12pt]{article}
\usepackage[utf8]{inputenc}
\usepackage{amsmath, amsthm, amssymb, amsfonts}
\usepackage{hyperref}

\newtheorem{theorem}{Theorem}[section]
\newtheorem{lemma}[theorem]{Lemma}
\newtheorem{definition}[theorem]{Definition}
\newtheorem{conjecture}[theorem]{Conjecture}

\title{Partial Proof Architecture for the Erd\H{o}s-Szemer\'edi Sum-Product Conjecture}
\author{Charles EDOU NZE\thanks{Charles EDOU NZE, chercheur indépendant}}
\date{\today}

\begin{document}

\maketitle

\begin{abstract}
This document outlines a rigorous partial resolution strategy for the Erd\H{o}s-Szemer\'edi Sum-Product Conjecture, utilizing incidences in discrete geometry and combinatorial bounds. We establish axiomatic foundations and step-by-step lemmas leading towards improved bounds for sum and product sets.
\end{abstract}

\section{Axiomatic Definitions}

We define the primary algebraic structures. Let $\mathbb{N}$ denote the set of natural numbers, and let $\mathbb{R}$ denote the set of real numbers.

\begin{definition}[Sum Set]
For a finite subset $A \subset \mathbb{R}$, the sum set is defined as:
$$ A + A = \{ x + y \mid x, y \in A \} $$
\end{definition}

\begin{definition}[Product Set]
For a finite subset $A \subset \mathbb{R}$, the product set is defined as:
$$ A \cdot A = \{ x \cdot y \mid x, y \in A \} $$
\end{definition}

\begin{conjecture}[Erd\H{o}s-Szemer\'edi, 1983]
For any $\varepsilon > 0$, there exists a constant $c(\varepsilon) > 0$ such that for any finite set $A \subset \mathbb{N}$,
$$ \max(|A + A|, |A \cdot A|) \geq c(\varepsilon) |A|^{2 - \varepsilon}. $$
\end{conjecture}

\section{Contextual Literature Research}
The foundational result by Elekes (1997) established a lower bound of $O(|A|^{5/4})$ using the Szemer\'edi-Trotter theorem on point-line incidences. Recent advancements by Solymosi, Konyagin, and Shkredov have improved this bound iteratively. The underlying methodology relies heavily on incidence geometry over the real plane $\mathbb{R}^2$. A profound analogy exists between this problem and the bounds on intersections of curves over finite fields, which are constrained by the polynomial method (as seen in Dvir's resolution of the finite field Kakeya problem).

\section{Lemmas and Proof Strategy}

We proceed by formalizing the incidence geometry approach.

\begin{lemma}[Point-Line Incidence Bound]
Let $P \subset \mathbb{R}^2$ be a finite set of points, and $L$ a finite set of lines in $\mathbb{R}^2$. The number of incidences $I(P, L) = |\{ (p, l) \in P \times L \mid p \in l \}|$ is bounded by:
$$ I(P, L) \leq c \left( |P|^{2/3}|L|^{2/3} + |P| + |L| \right) $$
for some absolute constant $c > 0$.
\end{lemma}

\begin{proof}[Proof of Lemma 1]
We construct a cell decomposition of $\mathbb{R}^2$ using polynomial partitioning. Let $r = |P|^{2/3}|L|^{-1/3}$. If $|L| > |P|^2$ or $|P| > |L|^2$, the trivial bounds apply. Otherwise, by the polynomial partitioning theorem, there exists a polynomial $f(x,y)$ of degree $O(\sqrt{r})$ whose zero set $Z(f)$ partitions $\mathbb{R}^2$ into $O(r)$ open cells, each containing at most $O(|P|/r)$ points.
By analyzing the incidences within the cells and on the algebraic variety $Z(f)$, and summing these contributions, we deduce the upper bound. Every line not fully contained in $Z(f)$ intersects it in at most $\deg(f) = O(\sqrt{r})$ points. Summing over all lines and cells yields the stated bound.
\end{proof}

\begin{lemma}[Elekes' Construction]
If $A \subset \mathbb{R}$, we define the point set $P = (A + A) \times (A \cdot A)$ and the line set $L = \{ y = a(x - b) \mid a, b \in A \}$. Then $I(P, L) \geq |A|^3$.
\end{lemma}

\begin{proof}[Proof of Lemma 2]
For each pair $(a, b) \in A \times A$, we define the line $l_{a,b}$ given by the equation $y = a(x - b)$. There are $|A|^2$ such lines, so $|L| = |A|^2$.
For any element $c \in A$, let $x = b + c \in A + A$. Then $y = a(x - b) = a(c) \in A \cdot A$.
The point $(x, y) = (b+c, ac)$ clearly belongs to $P = (A+A) \times (A \cdot A)$.
Thus, for a fixed line $l_{a,b}$, there are $|A|$ distinct choices of $c \in A$, yielding $|A|$ distinct points $(x,y) \in P$ on $l_{a,b}$.
Summing over all $|A|^2$ lines, the total number of incidences is at least $|A|^2 \cdot |A| = |A|^3$.
\end{proof}

\section{Partial Proof of the Sum-Product Bound}

By combining Lemma 1 and Lemma 2, we deduce a lower bound on $\max(|A+A|, |A \cdot A|)$.
Assume for contradiction that both $|A+A| \leq K|A|$ and $|A \cdot A| \leq K|A|$ for some small $K$.
Then $|P| = |A+A| \cdot |A \cdot A| \leq K^2 |A|^2$.
From Lemma 2, $I(P, L) \geq |A|^3$.
From Lemma 1,
$$ I(P, L) \leq c \left( |P|^{2/3}|L|^{2/3} + |P| + |L| \right) $$
Substituting $|L| = |A|^2$ and $|P| \leq K^2 |A|^2$:
$$ |A|^3 \leq c \left( (K^2 |A|^2)^{2/3}(|A|^2)^{2/3} + K^2 |A|^2 + |A|^2 \right) $$
$$ |A|^3 \leq c \left( K^{4/3} |A|^{4/3} |A|^{4/3} + K^2 |A|^2 + |A|^2 \right) $$
$$ |A|^3 \leq c \left( K^{4/3} |A|^{8/3} + K^2 |A|^2 + |A|^2 \right) $$
For large $|A|$, the dominant term is $K^{4/3} |A|^{8/3}$. Thus:
$$ |A|^3 \leq c K^{4/3} |A|^{8/3} $$
Dividing by $|A|^{8/3}$ gives $|A|^{1/3} \leq c K^{4/3}$.
Therefore, $K \geq c' |A|^{1/4}$.
This implies that $\max(|A+A|, |A \cdot A|) \geq c' |A|^{5/4}$.

\section{Autoformalization Architecture}
The following definitions and theorems are structured for future formalization in Lean 4.

\begin{verbatim}
import Mathlib.Data.Real.Basic
import Mathlib.Data.Finset.Basic

variable {A : Finset \mathbb{R}}

def sum_set (A : Finset \mathbb{R}) : Finset \mathbb{R} :=
  (A \times A).image (\lambda p => p.1 + p.2)

def prod_set (A : Finset \mathbb{R}) : Finset \mathbb{R} :=
  (A \times A).image (\lambda p => p.1 * p.2)

theorem erdos_szemeredi_weak_bound (A : Finset \mathbb{R}) :
  \exists c : \mathbb{R}, c > 0 \land (sum_set A).card * (prod_set A).card \ge c * (A.card ^ (5/2)) := by
  sorry
\end{verbatim}

\end{document}
